\documentclass[12pt]{article}

\usepackage[a4paper, margin=1.8cm]{geometry}
\usepackage{amsmath}
\usepackage{siunitx}
\usepackage{enumitem}

\setlength{\parindent}{0pt}

\begin{document}

\begin{center}
    \Large \textbf{GPE, KE -- Calculations and Conversions}
\end{center}

\[
E_p = mgh \quad\quad E_k = \frac{1}{2}mv^2 \quad\quad g = \SI{9.8}{m/s^2}
\]

\begin{enumerate}

\item A \SI{4}{kg} object is lifted \SI{3}{m} above the ground. Calculate its gravitational potential energy.

\item A \SI{2}{kg} object is moving at \SI{6}{m/s}. Calculate its kinetic energy.

\item A box gains \SI{490}{J} of gravitational potential energy when lifted \SI{5}{m}. Calculate the mass of the box.

\item A ball has \SI{160}{J} of kinetic energy and a mass of \SI{5}{kg}. Calculate its velocity.

\item A \SI{12}{kg} object has \SI{2352}{J} of gravitational potential energy. Calculate its height above the ground.

\item A roller coaster car of mass \SI{500}{kg} starts from rest at the top of a \SI{20}{m} hill. Calculate its gravitational potential energy at the top, then calculate its speed at the bottom assuming all GPE is converted into KE.

\item A \SI{0.5}{kg} tennis ball is thrown vertically upwards at \SI{12}{m/s}. Calculate its starting kinetic energy, then calculate the maximum height it reaches.

\item A \SI{2}{kg} rock is dropped from a bridge \SI{15}{m} above the water. Calculate its starting gravitational potential energy, then calculate its speed just before it hits the water.

\item A skydiver of mass \SI{75}{kg} falls \SI{100}{m}. Calculate the gravitational potential energy lost, then calculate the speed they would reach if all of this energy became kinetic energy.

\item A \SI{60}{kg} skateboarder rolls down a ramp from a height of \SI{2.5}{m}. Calculate the GPE at the top, then calculate the speed at the bottom.

\item A \SI{0.2}{kg} arrow is fired vertically upwards at \SI{50}{m/s}. Calculate its starting kinetic energy, then calculate the maximum height it could reach.

\item A \SI{1200}{kg} car rolls down a hill and loses \SI{58800}{J} of gravitational potential energy. Calculate the height lost, then calculate the speed gained if all the lost GPE becomes KE.

\item A \SI{10}{kg} sledge starts from rest at the top of a snowy slope \SI{8}{m} high. Calculate its GPE at the top, then calculate its speed at the bottom.

\item A \SI{0.8}{kg} drone rises vertically and gains \SI{156.8}{J} of gravitational potential energy. Calculate the height gained. If it then falls from this height, calculate its speed just before reaching the ground.

\item A \SI{70}{kg} diver jumps from a \SI{52}{m} platform. Calculate their GPE at the top, then calculate their impact speed ignoring air resistance.

\end{enumerate}

\newpage

\begin{center}
    \Large \textbf{Answer Key}
\end{center}

\begin{enumerate}

\item
\[
E_p = mgh = 4 \times 9.8 \times 3 = \SI{117.6}{J}
\]

\item
\[
E_k = \frac{1}{2}mv^2 = \frac{1}{2} \times 2 \times 6^2 = \SI{36}{J}
\]

\item
\[
m = \frac{E_p}{gh} = \frac{490}{9.8 \times 5} = \SI{10}{kg}
\]

\item
\[
v = \sqrt{\frac{2E_k}{m}} = \sqrt{\frac{2 \times 160}{5}} = \SI{8}{m/s}
\]

\item
\[
h = \frac{E_p}{mg} = \frac{2352}{12 \times 9.8} = \SI{20}{m}
\]

\item
\[
E_p = mgh = 500 \times 9.8 \times 20 = \SI{98000}{J}
\]
\[
E_k = \SI{98000}{J}
\]
\[
v = \sqrt{\frac{2E_k}{m}} = \sqrt{\frac{2 \times 98000}{500}} = \SI{19.8}{m/s}
\]

\item
\[
E_k = \frac{1}{2}mv^2 = \frac{1}{2} \times 0.5 \times 12^2 = \SI{36}{J}
\]
\[
h = \frac{E_p}{mg} = \frac{36}{0.5 \times 9.8} = \SI{7.35}{m}
\]

\item
\[
E_p = mgh = 2 \times 9.8 \times 15 = \SI{294}{J}
\]
\[
E_k = \SI{294}{J}
\]
\[
v = \sqrt{\frac{2E_k}{m}} = \sqrt{\frac{2 \times 294}{2}} = \SI{17.1}{m/s}
\]

\item
\[
E_p = mgh = 75 \times 9.8 \times 100 = \SI{73500}{J}
\]
\[
E_k = \SI{73500}{J}
\]
\[
v = \sqrt{\frac{2E_k}{m}} = \sqrt{\frac{2 \times 73500}{75}} = \SI{44.3}{m/s}
\]

\item
\[
E_p = mgh = 60 \times 9.8 \times 2.5 = \SI{1470}{J}
\]
\[
E_k = \SI{1470}{J}
\]
\[
v = \sqrt{\frac{2E_k}{m}} = \sqrt{\frac{2 \times 1470}{60}} = \SI{7}{m/s}
\]

\item
\[
E_k = \frac{1}{2}mv^2 = \frac{1}{2} \times 0.2 \times 50^2 = \SI{250}{J}
\]
\[
h = \frac{E_p}{mg} = \frac{250}{0.2 \times 9.8} = \SI{127.6}{m}
\]

\item
\[
h = \frac{E_p}{mg} = \frac{58800}{1200 \times 9.8} = \SI{5}{m}
\]
\[
v = \sqrt{\frac{2E_k}{m}} = \sqrt{\frac{2 \times 58800}{1200}} = \SI{9.9}{m/s}
\]

\item
\[
E_p = mgh = 10 \times 9.8 \times 8 = \SI{784}{J}
\]
\[
E_k = \SI{784}{J}
\]
\[
v = \sqrt{\frac{2E_k}{m}} = \sqrt{\frac{2 \times 784}{10}} = \SI{12.5}{m/s}
\]

\item
\[
h = \frac{E_p}{mg} = \frac{156.8}{0.8 \times 9.8} = \SI{20}{m}
\]
\[
E_k = \SI{156.8}{J}
\]
\[
v = \sqrt{\frac{2E_k}{m}} = \sqrt{\frac{2 \times 156.8}{0.8}} = \SI{19.8}{m/s}
\]

\item
\[
E_p = mgh = 70 \times 9.8 \times 52 = \SI{35672}{J}
\]
\[
E_k = \SI{35672}{J}
\]
\[
v = \sqrt{\frac{2E_k}{m}} = \sqrt{\frac{2 \times 35672}{70}} = \SI{31.9}{m/s}
\]

\end{enumerate}

\end{document}